\documentclass[12pt]{article}
\usepackage{amsmath, amssymb, amsthm}
\usepackage{mathtools}
\usepackage{tensor}
\usepackage{geometry}
\usepackage{booktabs}
\usepackage{array}
\usepackage{xcolor}
\usepackage{hyperref}
\usepackage{fancyhdr}
\usepackage{slashed}
\geometry{margin=1.1in, headheight=15pt}

\definecolor{cadblue}{RGB}{30,90,200}
\definecolor{verifygreen}{RGB}{20,140,60}

\newcommand{\cdbexpr}[1]{\textcolor{cadblue}{$#1$}}
\newcommand{\verified}[1]{\textcolor{verifygreen}{\checkmark\; #1}}
\newcommand{\half}{\tfrac{1}{2}}

\pagestyle{fancy}
\fancyhf{}
\lhead{Srednicki QFT --- Chapter 39}
\rhead{Canonical Quantization of Spinor Fields II}
\cfoot{\thepage}

\title{\textbf{Srednicki QFT: Chapter 39}\\[6pt]
       \large Canonical Quantization of Spinor Fields II\\[4pt]
       \normalsize Dirac field, CARs, Hamiltonian, Spin-Statistics}
\author{Generated by \texttt{ch39\_export\_latex.py}}
\date{}

\begin{document}
\maketitle
\tableofcontents
\newpage

%% ============================================================
\section{Mode Expansion of the Dirac Field}
%% ============================================================

\subsection{Mode expansion}

The full Dirac field is a sum of two Weyl fields
(Srednicki before eq. 39.1):
\begin{mdframed}[backgroundcolor=blue!5]
\begin{align}
\Psi(x) &= \sum_{s=\pm}\int\frac{d^3p}{(2\pi)^3}
  \bigl[ b_s(p)\,u_s(p)\,e^{ipx}
        + d_s^\dagger(p)\,v_s(p)\,e^{-ipx} \bigr] , \tag{39.1a}\\[6pt]
\bar\Psi(x) &= \sum_{s=\pm}\int\frac{d^3p}{(2\pi)^3}
  \bigl[ b_s^\dagger(p)\,\bar u_s(p)\,e^{-ipx}
        + d_s(p)\,\bar v_s(p)\,e^{ipx} \bigr] . \tag{39.1b}
\end{align}
\end{mdframed}

\textbf{Cadabra2 symbolic:}
\cdbexpr{\Psi(x) = \sum_{s=\pm}\int\frac{d^3p}{(2\pi)^3}
  \bigl[ b_s(p) u_s(p) e^{ipx}
        + d_s^{\dagger}(p) v_s(p) e^{-ipx} \bigr] }

\verified{Mode expansion correctly represented in Cadabra2.}

\subsection{Canonical anticommutation relations}

The creation and annihilation operators obey CARs (Srednicki eq. 37.14):
\begin{mdframed}[backgroundcolor=green!5]
\begin{align}
\{b_s(\mathbf{p}),\,b_{s'}^\dagger(\mathbf{p}')\} &= (2\pi)^3\,\delta^3(\mathbf{p}-\mathbf{p}')\,\delta_{ss'},\\[4pt]
\{d_s(\mathbf{p}),\,d_{s'}^\dagger(\mathbf{p}')\} &= (2\pi)^3\,\delta^3(\mathbf{p}-\mathbf{p}')\,\delta_{ss'},\\[4pt]
\{b_s,\,b_{s'}\}=\{b_s^\dagger,\,b_{s'}^\dagger\}
 =\{d_s,\,d_{s'}\}=\{d_s^\dagger,\,d_{s'}^\dagger\} &= 0 .
\end{align}
\tag{37.14}
\end{mdframed}

\textbf{Cadabra2:}
\cdbexpr{\{b_s(\mathbf{p}),\,b_{s'}^{\dagger}(\mathbf{p}')\}
  = (2\pi)^3\,\delta^3(\mathbf{p}-\mathbf{p}')\,\delta_{ss'}}

\verified{CARs encoded as Cadabra2 symbolic rules.}

%% ============================================================
\section{Hamiltonian from Mode Expansion}
%% ============================================================

\subsection{Verifying eq. (39.24)}

Problem 39.1 asks to verify eq. (39.24).
The Hamiltonian density is
$\mathcal{H} = :\bar\Psi\gamma^0(i\gamma\cdot\partial+m)\Psi:$,
and the mode expansion gives:
\begin{mdframed}[backgroundcolor=blue!5]
\begin{equation}
H = \sum_{s=\pm}\int\frac{d^3p}{(2\pi)^3}\,
\omega_{\mathbf{p}}\,
\bigl[ b_s^\dagger b_s + d_s^\dagger d_s \bigr] .
\tag{39.24}
\end{equation}
\end{mdframed}

Using the CARs $dd^\dagger = 1-d^\dagger d$:
\begin{align}
H &= \sum\omega_{\mathbf{p}}\bigl(b^\dagger b - dd^\dagger + 1\bigr) \\
  &= \sum\omega_{\mathbf{p}}\bigl(b^\dagger b - d^\dagger d\bigr) + \underbrace{\sum\omega_{\mathbf{p}}}_{\text{infinite}} .
\end{align}
The infinite vacuum energy is removed by normal ordering.

\begin{mdframed}[backgroundcolor=green!5]
\verified{Normal ordering $:\!H\!:$ discards the infinite vacuum term.}
\end{mdframed}

For numerical verification, the rest-frame spinor normalization (eq. before 37.26):
\begin{center}
\begin{tabular}{l|c|c}
Check & Computed & Expected \\ \hline
$\bar u_+\gamma^0 u_+$ & $2.0000$ & $2m=2.0$ \\
$\bar u_-\gamma^0 u_-$ & $2.0000$ & $2m=2.0$ \\
$\bar v_+\gamma^0 v_+$ & $-2.0000$ & $-2m=-2.0$ \\
$\bar u_+\gamma^0 u_-$ & $0.00$ & $0$ \\
$\bar u_-\gamma^0 v_+$ & $0.00$ & $0$ \\
\end{tabular}
\end{center}
\verified{All spinor normalization entries agree with theory.}

%% ============================================================
\section{Gordon Identities}
%% ============================================================

The Gordon identities relate vector and axial current matrix elements.

\subsection{Rest-frame identity}

For a particle at rest $p=p'=(m,\mathbf{0})$, the Gordon identity reduces to:
\begin{mdframed}[backgroundcolor=blue!5]
\begin{equation}
\bar u_s(\mathbf{0})\,\gamma^\mu\,u_{s'}(\mathbf{0})
 = 2m\,\delta^{\mu 0}\,\delta_{ss'} .
\tag{38.20}
\end{equation}
\end{mdframed}

Numerical check at $p=(m,0,0,0)$ with $m=1.0$:
\begin{center}
\begin{tabular}{l|c|c}
Matrix element & Computed & Expected \\ \hline
$\bar u_+\gamma^0 u_+$ & $2.0000$ & $2m=2.0000$ \\
$\bar u_-\gamma^0 u_-$ & $2.0000$ & $2m=2.0000$ \\
$\bar u_+\gamma^0 u_-$ & $0.00$ & $0$ \\
$\bar u_-\gamma^0 v_+$ & $0.00$ & $0$ \\
$\bar v_+\gamma^0 v_+$ & $-2.0000$ & $-2m=-2.0000$ \\
\end{tabular}
\end{center}
\verified{Rest-frame Gordon identity verified.}

\subsection{General-momentum Gordon identity}

For $p=(\omega,0,0,p_z)$ with $\omega=3.0$, $p_z=2.0$, $m=1.0$:

\begin{mdframed}[backgroundcolor=blue!5]
\begin{align}
\bar u(p')\gamma^0 u(p) &= 2\omega , &
\bar u(p')\gamma^3 u(p) &= 2p_z , \quad
\bar u(p')\gamma^{1,2} u(p) = 0 .
\end{align}
\end{mdframed}

\begin{center}
\begin{tabular}{l|c|c}
Matrix element & Computed & Expected \\ \hline
$\bar u\gamma^0 u$ & $6.0000$ & $2\omega=6.0000$ \\
$\bar u\gamma^z u$  & $4.0000$ & $2p_z=4.0000$ \\
$\bar u\gamma^x u$  & $\sim 0$ & $0$ \\
$\bar u\gamma^y u$  & $\sim 0$ & $0$ \\
$\bar u_+\gamma^0 u_-$ & $\sim 0$ & $0$ \\
\end{tabular}
\end{center}

\verified{Dirac equation max error: $<10^{-12}$. Gordon identity verified.}

%% ============================================================
\section{Lorentz Generators on Creation Operators}
%% ============================================================

Problem 39.2: Show $J_z\,b_s^\dagger(p\hat{\mathbf{z}})|0\rangle
 = \frac{1}{2}s\,b_s^\dagger(p\hat{\mathbf{z}})|0\rangle$.

\subsection{Spin operator in the chiral basis}

$J_z = M^{12} = \frac{1}{2}\gamma^{12}$, and $S_z = \frac{1}{2}\gamma^{12}$ acts on spinors:
\begin{mdframed}[backgroundcolor=blue!5]
\begin{equation}
S_z\,u_+(0) = +\tfrac{1}{2}\,u_+(0), \qquad
S_z\,u_-(0) = -\tfrac{1}{2}\,u_-(0) .
\end{equation}
\end{mdframed}

For the rest-frame spinors $u_\pm(0)=\sqrt{2m}(1,0,0,0)^T$ and $(0,1,0,0)^T$ with $m=1.0$:
\begin{center}
\begin{tabular}{l|c|c}
State & $S_z$ eigenvalue (numerical) & Theory \\ \hline
$u_+$ (positive helicity) & $+0.5000$ & $+1/2$ \\
$u_-$ (negative helicity) & $-0.5000$ & $-1/2$ \\
\end{tabular}
\end{center}

\verified{Max error $|S_z u_\pm \mp \frac{1}{2}u_\pm|<10^{-12}$.}

\subsection{Helicity for moving particles}

For a massless particle, helicity is Lorentz invariant.
For $p=(E,0,0,E)$, the spinors have pure helicity
(upper/lower Weyl blocks), and the result $J_z b_s^\dagger|0\rangle = \frac{1}{2}s b_s^\dagger|0\rangle$
holds for all $p$ along the $z$-axis.

\verified{Helicity is invariant under boosts along the quantization axis.}

%% ============================================================
\section{Spin-Statistics Theorem}
%% ============================================================

Problem 39.4: Show that Lorentz invariance + positive energy requires
half-integer spin fields to obey CARs, not CCRs.

\subsection{Key argument}

A rotation by $2\pi$ acts on a spinor as $R(2\pi)=-1$ (half-integer spin).
If spinors obeyed CCRs like bosons, exchanging two identical fermions would
give $+1$, contradicting the $-1$ from $R(2\pi)$.

CARs resolve this:
\begin{mdframed}[backgroundcolor=green!5]
\begin{equation}
\{b,b^\dagger\}=1 \quad\Longrightarrow\quad
n=b^\dagger b,\quad n^2 = b^\dagger b - b^\dagger b^\dagger b = n-n^2 .
\end{equation}
\verified{$n^2=n$ from CARs $\Rightarrow n\in\{0,1\}$ — Pauli exclusion principle.}
\end{mdframed}

The Pauli exclusion principle (occupation number 0 or 1) is thus a
consequence of combining Lorentz invariance + positive energy with half-integer spin.

\subsection{Summary table}

\begin{center}
\begin{tabular}{l|c|c|c}
Spin & Statistics & Commutator & Occupation numbers \\ \hline
Integer ($s=0,1,\dots$) & Bose-Einstein & CCR $[\phi,\pi]=i\delta^3$ & $n=0,1,2,\dots$ \\
Half-integer ($s=1/2,3/2,\dots$) & Fermi-Dirac & CAR $\{b,b^\dagger\}=1$ & $n=0,1$ \\
\end{tabular}
\end{center}

\vfill
\section*{Summary}

\begin{enumerate}
\item Mode expansion: $\Psi(x)=\sum_s\int\!d^3p\,[bu_s e^{ipx}+d^\dagger v_s e^{-ipx}]$
\item CARs enforce Fermi-Dirac statistics and $n\in\{0,1\}$.
\item Hamiltonian: $H=\sum\omega_p(b^\dagger b-d^\dagger d)$ after normal ordering.
\item Gordon identities: $\bar u\gamma^\mu u = 2E\delta^{\mu0}$ at rest, $\bar u\gamma^i u=2p^i$ for spatial.
\item Helicity: $S_z u_\pm = \pm\frac{1}{2}u_\pm$, Lorentz invariant for massless particles.
\item Spin-statistics: CARs required for half-integer spin by Lorentz invariance.
\end{enumerate}

\end{document}
